\chapter{Transferring a recipe to wider models}
\label{ch:width}
\source{P1:67--73}

Retuning a large model by an exhaustive search is expensive. The goal is to discover a rule on a cheap model and transfer it to a wider one.

\section{The empirical problem with standard parameterization}
Keep four transformer layers and vary the embedding width $d$. With standard parameterization, the best learning rate decreases as width grows. Since the dominant parameter count scales approximately as $Ld^2$, a doubling of width roughly quadruples parameters at fixed depth.

\figpair{width_sweep.png}{width_fit.png}{P1:67--68: wider models require smaller learning rates in this standard-parameterization experiment. The fitted exponent printed on P1:68 is about $-1.58$.}{width-standard}{}

This small sweep supports a local empirical trend, not a precise universal law. P1:68 gives a fitted exponent near $-1.58$, while P1:70 gives a value near $-1.34$. Both support the qualitative trend. They do not justify reporting a uniquely measured exponent without qualification.

\section{Why width changes the meaning of a parameter update}
Consider one component of a linear layer:
\[
y_i=\sum_{j=1}^{d}W_{ij}x_j,
\qquad
\Delta y_i=\sum_{j=1}^{d}\Delta W_{ij}x_j.
\]
An Adam update can have a coordinate scale of order $\eta$, because its numerator is divided by an estimate of its root-mean-square magnitude. If update contributions add coherently, the resulting feature change can scale like $d\eta$. Keeping a feature change of order one then suggests $\eta\propto d^{-1}$.

This argument provides intuition rather than a complete proof. If terms were independent and centered, their sum would scale differently, like $\sqrt d$ in standard deviation. Learned updates are correlated with activations and gradients, so independence is generally the wrong assumption. A valid theory must track those correlations, initialization scales, forward activations, backward signals, and the output layer together.

\section{What maximal update parameterization proposes}
The Tensor Programs V paper~\cite{mup} develops \emph{maximal update parameterization}, written $\mu$P, to preserve nontrivial feature learning as width grows. Its practical payoff is often called \emph{$\mu$Transfer}: tune a base learning rate at small width, then reuse it with prescribed parameter-group scaling at larger width.

For width-only scaling with Adam, the operational intuition is:
\begin{itemize}
\item scale learning rates for hidden weight matrices inversely with the width multiplier;
\item treat vector-like parameters, such as normalization gains, differently;
\item use the appropriate embedding and readout conventions;
\item keep initialization and forward-pass scaling consistent with the chosen parameterization.
\end{itemize}
Rescale weight decay for matrices to preserve the decay timescale. If a particular parameter group's learning rate becomes $\eta/r$, choosing its decay coefficient as $r\lambda$ preserves the product $\eta\lambda$. This is a statement about an AdamW implementation that multiplies decay by learning rate; optimizer conventions must be checked.

\extension{The rule ``divide every learning rate by width'' is not $\mu$P. Different tensors have different width-dependent dimensions. The reference implementation~\cite{mupcode} records base and target shapes and provides a compatible readout layer, initialization, and optimizer handling. A correct implementation needs that full specification. The simplified recipe above is therefore specific to this model family.}

\figstack{mup_paper.png}{mup_coordinates.png}{P1:71--72: the transfer idea and a coordinate check. Coordinate checks test whether activation scales remain controlled as width changes; they do not by themselves establish optimal validation loss.}{mup-theory}{}

\section{How to validate transfer experimentally}
Before an expensive run, compare several widths using the same \emph{base} learning-rate grid. Under successful transfer, the curves can move downward as model capacity improves while their minima remain near the same horizontal coordinate.

\fig{mup_sweeps.png}{P1:73: chess experiments after width-dependent rescaling. The optimal base learning rates approximately align, making a small-model sweep useful for wider models.}{mup-chess}

Two different checks are useful. A \emph{coordinate check} measures forward activations or feature changes across widths for a small number of steps. It catches scaling errors before a long run. A \emph{transfer check} measures final loss across learning rates and widths. It tests the desired practical result. Passing one does not imply passing the other.

Keep the scope of the claim narrow. Changing depth, attention conventions, optimizer, batch size, sequence length, or data distribution can require additional rules or a new calibration. In particular, $\mu$P addresses width dependence; the horizon dependence established in the preceding chapter remains a separate issue.

\takeaway{Learning rates depend on parameterization. Transferring small-model tuning to larger models requires compatible width scaling, initialization, readouts, and parameter groups.}
\FloatBarrier
